\documentclass[11pt,a4paper]{article}
\usepackage[utf8]{inputenc}
\usepackage[UTF8]{ctex} 
\usepackage{amsmath, amssymb, amsfonts, amsthm}
\usepackage{geometry}
\usepackage{float}
\usepackage{bm}
\usepackage{tcolorbox}
\usepackage[colorlinks=true, linkcolor=blue, unicode=true]{hyperref}

\geometry{left=2cm, right=2cm, top=2.5cm, bottom=2.5cm}

\title{基于 DBI 指标的分层初始化 K-means \\ (DHIKM: DBI Based Hierarchical Initialization)}
\date{\today}

\newtheorem{property}{性质}
\newtheorem{definition}{定义}

\begin{document}
\maketitle

\section{研究背景与改进分析}
传统的 K-means 及其分层初始化方法（HI）存在一个共同缺陷：需要预先指定聚类数目 $K$。本文提出了 \textbf{DHIKM} 算法，通过集成 Davies Bouldin Index (DBI)，在分层采样的顶部自动确定最优的 $K$。

\section{数学定义：Davies Bouldin Index (DBI)}
DBI 是衡量聚类效果的内部指标，其值越小表示聚类效果越好（类内紧凑，类间分离）。

\begin{definition}[类内分散度 $S_i$]
对于第 $i$ 个簇，其分散度定义为类内各点到质心 $Z_i$ 的平均距离：
\begin{equation}
    S_i = \left( \frac{1}{T_i} \sum_{j=1}^{T_i} |X_j - Z_i|^p \right)^{1/p}
\end{equation}
（通常取 $p=2$ 即欧氏距离）。
\end{definition}

\begin{definition}[类间距离 $M_{i,j}$]
第 $i$ 簇与第 $j$ 簇质心之间的距离：
\begin{equation}
    M_{i,j} = ||Z_i - Z_j||_p
\end{equation}
\end{definition}

\begin{definition}[DBI 指标]
\begin{equation}
    DBI = \frac{1}{K} \sum_{i=1}^{K} R_i, \quad \text{其中 } R_i = \max_{j \neq i} \left( \frac{S_i + S_j}{M_{i,j}} \right)
\end{equation}
\end{definition}

\section{DHIKM 算法流程}



\begin{tcolorbox}[colback=green!5, colframe=green!40!black, title=DHIKM 算法核心步骤]
\begin{enumerate}
    \item \textbf{自底向上采样}：同前作，将原始数据集 $L_0$ 逐层采样至最顶层 $L_m$（数据量足够小）。
    \item \textbf{自动寻找 $K$}：
    \begin{itemize}
        \item 在 $L_m$ 层，设定 $K$ 的搜索范围 $[K_{min}, K_{max}]$。
        \item 对每个 $K$，选取权值最大的 $K$ 个点作为初始中心进行加权 K-means。
        \item 计算每个 $K$ 值对应的 $DBI$。
        \item 选取使 $DBI$ 最小的 $K$ 作为最优聚类数。
    \end{itemize}
    \item \textbf{投影与精炼}：使用最优 $K$ 及对应的初始中心，逐层向下投影并执行 K-means。
\end{enumerate}
\end{tcolorbox}

\section{性质 3 的证明：DBI 在层间的一致性}
\textbf{性质 3}：若在第 $k+1$ 层获得了较小的 DBI 值，则该 $K$ 值在第 $k$ 层也极大概率是有效的。

\begin{proof}
根据前作性质 2 的结论，$2 \cdot X_{k+1} - X_k \le 1$。
\begin{itemize}
    \item \textbf{类内分散度 $S$ 的关系}：由于坐标投影误差 $\le 1$，第 $k$ 层的分散度 $S^k \approx 2 \cdot S^{k+1}$。
    \item \textbf{类间距离 $M$ 的关系}：同理，$M^k \approx 2 \cdot M^{k+1}$。
    \item \textbf{比例不变性}：在计算 $R_i = (S_i + S_j)/M_{i,j}$ 时，分子分母的倍数关系抵消。
\end{itemize}
因此，在极小规模的顶层（计算代价极低）确定的最优 $K$ 可以直接应用于原始大规模数据层。
\end{proof}

\section{实验与性能结论}
\begin{itemize}
    \item \textbf{计算效率}：在顶层（数据量约为原始的 $1/4^m$）搜索 $K$ 值，比在原始层搜索速度快数百倍。
    \item \textbf{准确性}：在 UCI 数据集（如 Iris, Wine, Seeds）实验中，DHIKM 准确识别了真实的聚类数目。
    \item \textbf{结论}：该方法解决了 K-means 初始化中两个最棘手的问题：初始中心的选择和聚类数目的确定。
\end{itemize}

\end{document}